% Two-page evidence-linked report on the nonisothermal CSTR.
%
% Build:
%   latexmk -pdf report.tex          (from the report/ directory)
%   latexmk -C                       (clean)
%
% Three rules the claim audit in Workshop 2 will check:
%
%   1. Every numerical value traces to a file in ../results/.
%   2. Every figure names the command that generated it.
%   3. Every citation is a source you opened yourself.
%
% The \artifact and \generatedby macros below exist so that tracing is
% mechanical rather than a memory exercise. Use them.

\documentclass[11pt]{article}
\usepackage[margin=1in]{geometry}
\usepackage{amsmath}
\usepackage{graphicx}
\usepackage{booktabs}
\usepackage[hidelinks]{hyperref}

% Record where a number or a figure came from. Redefine to \marginpar, a
% footnote, or nothing at all -- but keep the argument honest.
\newcommand{\artifact}[1]{\textsuperscript{[\texttt{#1}]}}
\newcommand{\generatedby}[1]{\par\smallskip\footnotesize Generated by
  \texttt{#1}.\normalsize}

\title{Steady-state multiplicity in a jacketed nonisothermal CSTR}
\author{Your name}
\date{\today}

\begin{document}
\maketitle

\section{Research question}
% One or two sentences. What did you actually ask? Not what the field asks.
% A question your computation can answer is better than an important one it
% cannot.

\section{Model and assumptions}
% State the balances you implemented, with units and sign conventions.
%
% Every assumption your CODE makes belongs here, including the ones that are
% easy to forget because they are invisible in the equations -- constant coolant
% temperature, constant volume, constant physical properties, a single
% irreversible reaction. An assumption used in the code and omitted from the
% report is a real error, not a stylistic one.

At steady state,
\begin{align}
  0 &= q\,(C_{Af} - C_A) - V\,k(T)\,C_A, \\
  0 &= q\,\rho\,C_p\,(T_f - T) + (-\Delta H_{\mathrm{rxn}})\,V\,k(T)\,C_A
       - UA\,(T - T_c),
\end{align}
with $k(T) = k_0 \exp(-E/RT)$ and $\Delta H_{\mathrm{rxn}} < 0$.

\section{Computational method}
% Solver, initial-guess strategy, how you decided two roots were distinct, and
% the tolerances. State the software versions or point at the baseline artifact
% that records them\artifact{results/baseline.json}.

\section{Principal result}
% The numbers. Each one carries an \artifact{} pointing at the file it came
% from. Round consistently, and round the same value the same way everywhere.

\begin{figure}[htbp]
  \centering
  % \includegraphics[width=0.7\textwidth]{figures/steady_state_locus.png}
  \caption{Steady-state reactor temperature against coolant temperature.}
  \label{fig:locus}
  \generatedby{python scripts/reproduce.py}
\end{figure}

\section{Comparison with literature}
% For each comparison, say which of these it is -- and say it in the text, not
% only in your head:
%
%   numerical      same quantity, comparable conditions, values agree within a
%                  stated tolerance
%   qualitative    same phenomenon, different conditions or parameters
%   not comparable different system, assumptions, or parameters
%   insufficiently specified  the source does not report what you would need
%
% Different parameters are not a contradiction. They are also not agreement.

\section{Limitations}
% The honest list. At minimum: what your computation does NOT establish.
%
% If you want to say something about stability, check first whether you computed
% it. A locus of steady states of the algebraic balances is not a stability
% result, however suggestive its shape.

\section{Reproducibility}
% The command that regenerates every number and figure above, the commit these
% results came from, and where the environment is specified.

\bibliographystyle{plain}
\bibliography{ref}

\end{document}
